\begin{figure}[H]
\centering
\begin{tikzpicture}[
  node distance=10mm and 13mm,
  state/.style={draw=VIULinkBlue, rounded corners=1.5mm, fill=VIULight,
    minimum width=34mm, minimum height=13mm, align=center, font=\small},
  dual/.style={state, draw=VIUOrange, fill=VIUOrange!10},
  flow/.style={->, >=stealth, thick, draw=VIUDeep}
]
\node[state] (primal) {Primal\\$x_{iks}$};
\node[dual, right=of primal] (wrench) {Precio wrench\\$\lambda_k=r_k$};
\node[dual, below=of wrench] (slot) {Precio de slot\\$\mu_{ks}\geq0$};
\node[state, left=of slot] (equilibrium) {v-GNE / KKT\\relajación convexa};

\draw[flow] (primal) -- node[above, font=\scriptsize] {$\widetilde w-Hx$} (wrench);
\draw[flow] (wrench) -- node[right, font=\scriptsize] {$H^\top\lambda$} (slot);
\draw[flow] (slot) -- node[below, font=\scriptsize] {complementariedad} (equilibrium);
\draw[flow] (equilibrium) -- node[left, font=\scriptsize] {proyección} (primal);
\draw[flow, dashed] (primal) -- node[below, font=\scriptsize] {$Bx-\mathbf1$} (slot);
\end{tikzpicture}
\caption[Relación entre primal--dual, KKT y v-GNE]{El residual wrench actúa como precio vectorial; los multiplicadores de slot imponen la restricción compartida. La equivalencia se limita a la relajación convexa.}
\viuownsource
\label{fig:sp3-kkt-game}
\end{figure}
